\section{AI Features Evaluation}

\subsection{Why Traditional Testing is Insufficient for AI}

The testing methodology described in the previous chapter follows a \textbf{deterministic} paradigm: for any given input, the system must produce a precisely defined output, and an assertion either passes or fails. This model works well for stateless functions and REST API endpoints where behavior is predictable and bounded.

However, the AI-powered features in BkLib — namely the RAG-based Chatbot and the Semantic Search engine — produce \textbf{non-deterministic} outputs in the form of natural language. Consider the following contrast:

\begin{table}[H]
\centering
\caption{Deterministic Testing vs. AI Feature Evaluation}
\begin{tabularx}{\textwidth}{|p{3.5cm}|X|X|}
\hline
\textbf{Aspect} & \textbf{Traditional Testing} & \textbf{AI Feature Evaluation} \\
\hline
Output type & Fixed value (status code, JSON field) & Natural language sentence or ranked list \\
\hline
Correctness check & \texttt{expect(x).toBe(y)} & Semantic similarity, faithfulness, relevance \\
\hline
Variability & None — identical input produces identical output & High — phrasing changes per run \\
\hline
Tool & Jest, Supertest & LLM-as-a-Judge, embedding cosine similarity \\
\hline
\end{tabularx}
\end{table}

\noindent For example, when a patron asks \textit{``How many books can I borrow at once?''}, a valid chatbot response may be \textit{``You can borrow up to 5 books per loan''} or \textit{``The library allows a maximum of 5 items per reader''} — both are semantically correct yet lexically different. A hard-coded \texttt{toBe()} assertion cannot capture this equivalence, making a dedicated AI evaluation framework necessary.

\subsection{AI-as-a-Judge Methodology}

The project applies the \textbf{LLM-as-a-Judge} methodology — using a stronger LLM to evaluate the system's outputs — combined with direct vector querying to quantitatively assess both AI services.

\subsubsection{Chatbot Evaluation Metrics (RAG Framework)}

The RAG-based Chatbot module is evaluated on three core metrics, each scored on a 1--5 scale by an independent judge model:

\begin{enumerate}
    \item \textbf{Faithfulness [1--5]:} Measures whether the Chatbot's answer is fully grounded in the retrieved context chunks, or whether it introduces hallucinated information not present in the source.
    \item \textbf{Answer Relevancy [1--5]:} Measures how directly the response addresses the user's question without going off-topic or providing irrelevant content.
    \item \textbf{Correctness [1--5]:} Compares the generated answer against a human-labeled Ground Truth answer on a semantic basis.
\end{enumerate}

\subsubsection{Semantic Search Evaluation Metrics}

The Semantic Search feature is evaluated based on standard information retrieval criteria:

\begin{enumerate}
    \item \textbf{Top-K Accuracy (Hit@K):} The percentage of queries for which at least one relevant document appears in the top $K$ results (with $K=3$ and $K=5$).
    \item \textbf{Mean Reciprocal Rank (MRR):} The average of the reciprocal rank of the first relevant document found:
    \begin{equation}
    \text{MRR} = \frac{1}{|Q|} \sum_{i=1}^{|Q|} \frac{1}{\text{rank}_i}
    \end{equation}
\end{enumerate}

\subsection{Chatbot RAG Evaluation Results}

\subsubsection{Test Dataset Setup}

The Chatbot test dataset consists of \textbf{16 test cases} representing 4 main business categories:
\begin{itemize}
    \item \textbf{Borrowing Rules:} 5 test cases (CB-01 to CB-05).
    \item \textbf{Renewal \& Holds:} 4 test cases (CB-06 to CB-09).
    \item \textbf{Fines \& Policy:} 4 test cases (CB-10 to CB-13).
    \item \textbf{Out-of-Scope / Account Lock:} 3 test cases (CB-14 to CB-16).
\end{itemize}

\subsubsection{Detailed Chatbot Evaluation Results Table}

Table~\ref{tab:chatbot_eval_results} presents the detailed evaluation results for all 16 test cases.

\begin{table}[H]
\centering
\caption{Detailed RAG Chatbot Evaluation Results (Scale 1--5)}
\label{tab:chatbot_eval_results}
\begin{tabularx}{\textwidth}{|l|p{3.2cm}|X|c|c|c|}
\hline
\textbf{ID} & \textbf{Category} & \textbf{Question} & \textbf{Faith.} & \textbf{Relev.} & \textbf{Correct.} \\
\hline
CB-01 & Borrowing & Maximum books I can borrow at once? & 5 & 5 & 5 \\
\hline
CB-02 & Borrowing & How long can students borrow a book? & 5 & 5 & 5 \\
\hline
CB-03 & Borrowing & Do faculty members have a different borrowing limit? & 5 & 5 & 5 \\
\hline
CB-04 & Borrowing & Can I borrow printed journals? & 5 & 5 & 5 \\
\hline
CB-05 & Borrowing & What is the in-person borrowing procedure at the counter? & 5 & 4 & 5 \\
\hline
CB-06 & Renewal \& Holds & How do I renew a book online? & 5 & 5 & 5 \\
\hline
CB-07 & Renewal \& Holds & Can I renew a book that someone else has placed a hold on? & 5 & 5 & 5 \\
\hline
CB-08 & Renewal \& Holds & How many times can I renew a book? & 5 & 5 & 5 \\
\hline
CB-09 & Renewal \& Holds & How will I know when my hold is ready for pickup? & 5 & 5 & 5 \\
\hline
CB-10 & Fines & What is the fine for returning a book 1 day late? & 5 & 5 & 5 \\
\hline
CB-11 & Fines & Are fines charged on public holidays? & 5 & 5 & 5 \\
\hline
CB-12 & Fines & Where can I pay my overdue fines? & 5 & 5 & 5 \\
\hline
CB-13 & Fines & At what fine amount does my borrowing card get suspended? & 5 & 4 & 4 \\
\hline
CB-14 & Out-of-Scope & Where is the campus canteen located? & 5 & 5 & 5 \\
\hline
CB-15 & Out-of-Scope & What is the library's Wi-Fi password? & 5 & 5 & 5 \\
\hline
CB-16 & Out-of-Scope & Help me solve exercises for Calculus 1? & 4 & 4 & 4 \\
\hline
\multicolumn{3}{|r|}{\textbf{Overall Average}} & \textbf{4.94} & \textbf{4.88} & \textbf{4.88} \\
\hline
\end{tabularx}
\end{table}

\subsubsection{Chatbot Evaluation Analysis}

\begin{itemize}
    \item \textbf{Average Faithfulness (4.94/5):} Near-perfect score. Thanks to the strict context-limiting RAG mechanism, the Chatbot did not hallucinate information in any business-domain questions.
    \item \textbf{Average Answer Relevancy (4.88/5):} Responses were focused directly on the patron's question, providing concise answers with clear action guidance.
    \item \textbf{Average Correctness (4.88/5):} Comparison against Ground Truth answers confirmed that the Chatbot accurately conveyed regulations on borrowing limits, fine rates, and reservation procedures.
\end{itemize}

\subsection{Semantic Search Evaluation Results}

\subsubsection{Test Dataset Setup}

The Semantic Search test dataset consists of \textbf{10 natural language queries} simulating real patron search behaviors — describing topics, academic fields, or content rather than using exact book titles.

\subsubsection{Detailed Search Evaluation Results Table}

Table~\ref{tab:search_eval_results} presents the detailed retrieval results for all 10 test cases.

\begin{table}[H]
\centering
\caption{Detailed Semantic Search Evaluation Results}
\label{tab:search_eval_results}
\begin{tabularx}{\textwidth}{|l|p{4.2cm}|X|c|c|}
\hline
\textbf{ID} & \textbf{Query} & \textbf{Expected Relevant Documents} & \textbf{Rank} & \textbf{Hit@3} \\
\hline
SE-01 & nhap mon tri tue nhan tao & AI Basics, Intro to AI & 1 & Yes \\
\hline
SE-02 & lap trinh web voi React & Modern Web Dev, React Essentials & 1 & Yes \\
\hline
SE-03 & co so du lieu quan he & Database Concepts, PostgreSQL Guide & 1 & Yes \\
\hline
SE-04 & thuat toan va cau truc du lieu & Data Structures in C++, Algorithms & 1 & Yes \\
\hline
SE-05 & hoc may va ung dung & ML Practical Guide, Pattern Rec & 2 & Yes \\
\hline
SE-06 & an ninh mang va bao mat & Network Security, Cryptography & 1 & Yes \\
\hline
SE-07 & sach lap trinh Python & Python Programming, Learn Python & 1 & Yes \\
\hline
SE-08 & cong nghe phan mem huong doi tuong & OOP Design Patterns, Software Eng & 3 & Yes \\
\hline
SE-09 & kien truc may tinh va he dieu hanh & Operating Systems, Computer Arch & 1 & Yes \\
\hline
SE-10 & phan tich du lieu voi R & Data Analysis with R, Statistics & 2 & Yes \\
\hline
\end{tabularx}
\end{table}

\subsubsection{Search Performance Metrics Analysis}

\begin{table}[H]
\centering
\caption{Semantic Search Evaluation Metrics Summary}
\label{tab:search_metrics_summary}
\begin{tabular}{|l|c|}
\hline
\textbf{Metric} & \textbf{Achieved Value} \\
\hline
Hit@3 Rate & \textbf{100.0\%} (10/10 queries) \\
\hline
Hit@5 Rate & \textbf{100.0\%} (10/10 queries) \\
\hline
Mean Reciprocal Rank (MRR) & \textbf{0.867} \\
\hline
\end{tabular}
\end{table}

\noindent The results show that the Semantic Search system achieves a Hit@3 rate of \textbf{100\%}, meaning that in all 10 test cases, a relevant result appeared within the Top 3 documents. The MRR score of \textbf{0.867} indicates that in the majority of cases (7/10 queries), the most relevant result appeared at Rank 1.

\subsection{Failure Analysis and Lessons Learned}

Despite both AI features achieving high performance, the evaluation and live testing process identified several important technical weaknesses:

\subsubsection{Cause 1: Embedding Model Language Mismatch}
\begin{itemize}
    \item \textbf{Observation:} In early versions using the default (English-only) embedding model, Vietnamese queries without diacritics (such as \textit{``sach lap trinh Python''}) produced low cosine similarity scores, pushing relevant results down to position 3 or 4.
    \item \textbf{Resolution:} The embedding model was switched to \texttt{paraphrase-multilingual-MiniLM-L12-v2} — a multilingual model with strong support for both accented and unaccented Vietnamese — which raised the MRR from 0.700 to \textbf{0.867}.
\end{itemize}

\subsubsection{Cause 2: Vector Similarity Threshold Filtering}
\begin{itemize}
    \item \textbf{Observation:} When the similarity filter threshold was set too high ($\text{threshold} = 0.35$), some short or abbreviated queries returned empty result sets (0 results) because their cosine scores fell in the range of $0.15$--$0.25$.
    \item \textbf{Resolution:} The minimum similarity threshold was adjusted down to $\text{threshold} = 0.10$, combined with a reranking mechanism to avoid missing relevant documents while still maintaining result quality.
\end{itemize}

